\chapter{Conception et Réalisation}

\section{Introduction}

Ce chapitre documente les trois sprints de développement à travers lesquels
VM Marketplace a été construite. Chaque section de sprint s'ouvre sur son objectif et
son backlog, détaille les décisions de conception et les artefacts d'implémentation
clés, et se clôt par les difficultés rencontrées et la façon dont elles ont été
résolues.

% =============================================================================
\section{Sprint 1 --- Authentification et Provisionnement de VM}
% =============================================================================

\subsection{Objectif du Sprint}

Livrer un flux complet de bout en bout, de l'inscription de l'utilisateur à une VM
provisionnée et opérationnelle sur HCS, incluant le paiement Stripe et la diffusion
SSE de progression en temps réel.

\subsection{Backlog du Sprint}

\begin{table}[H]
\centering
\caption{Backlog Sprint 1}
\label{tab:sprint1_backlog}
\begin{tabularx}{\linewidth}{llX}
\toprule
\textbf{ID} & \textbf{Statut} & \textbf{Description} \\
\midrule
US-01 & Terminé & S'inscrire et se connecter avec e-mail / mot de passe \\
US-02 & Terminé & Configurer une VM via l'assistant multi-étapes \\
US-03 & Terminé & Payer via le formulaire Stripe \\
US-04 & Terminé & Provisionner la VM sur HCS via Terraform \\
US-05 & Terminé & Diffuser la progression du provisionnement via SSE \\
US-06 & Terminé & Envoyer un e-mail de confirmation après provisionnement \\
US-07 & Terminé & Consulter et supprimer ses VM dans le tableau de bord \\
\bottomrule
\end{tabularx}
\end{table}

\subsection{Conception de l'Authentification}

\paragraph{Hachage du mot de passe.} Les mots de passe des utilisateurs sont hachés
avec bcrypt (10 rounds) via la bibliothèque \texttt{passlib}. Le hash est stocké dans
la table \texttt{utilisateurs} ; le mot de passe en clair n'est jamais persisté.

\paragraph{Jetons JWT.} À la connexion réussie, le backend génère un JWT signé
contenant l'identifiant (\texttt{sub}) et la date d'expiration (\texttt{exp} = heure
actuelle + 30 minutes), en utilisant l'algorithme HS256. Le frontend stocke le jeton
dans le \texttt{localStorage} et l'attache en tant qu'en-tête \texttt{Bearer} à toutes
les requêtes API ultérieures.

La figure~\ref{fig:auth_seq} présente le diagramme de séquence d'authentification.

\begin{figure}[H]
\centering
\begin{tikzpicture}[
  box/.style={rectangle, draw, minimum width=2.8cm, minimum height=0.7cm,
              align=center, font=\small},
  arr/.style={-{Stealth}, thick},
]
  \node[box] (browser) at (0,0)    {Navigateur};
  \node[box] (api)     at (6,0)    {FastAPI};
  \node[box] (db)      at (11,0)   {SQLite};

  % Lignes de vie
  \draw[dashed] (browser.south) -- (0,-7);
  \draw[dashed] (api.south)     -- (6,-7);
  \draw[dashed] (db.south)      -- (11,-7);

  % Inscription
  \draw[arr] (0,-1) -- node[above,font=\scriptsize]{POST /api/register \{email, mdp\}} (6,-1);
  \draw[arr] (6,-1.6) -- node[above,font=\scriptsize]{INSERT utilisateur (hash bcrypt)} (11,-1.6);
  \draw[arr] (11,-2.2) -- node[above,font=\scriptsize]{200 OK} (6,-2.2);
  \draw[arr] (6,-2.8) -- node[above,font=\scriptsize]{201 \{message: "inscrit"\}} (0,-2.8);

  % Connexion
  \draw[arr] (0,-3.8) -- node[above,font=\scriptsize]{POST /api/login \{email, mdp\}} (6,-3.8);
  \draw[arr] (6,-4.4) -- node[above,font=\scriptsize]{SELECT utilisateur WHERE email=...} (11,-4.4);
  \draw[arr] (11,-5.0) -- node[above,font=\scriptsize]{Enregistrement utilisateur} (6,-5.0);
  \draw[arr] (6,-5.6) -- node[above,font=\scriptsize]{bcrypt.verify + signature JWT} (6,-5.6);
  \draw[arr] (6,-6.2) -- node[above,font=\scriptsize]{200 \{access\_token, token\_type\}} (0,-6.2);

  \node[font=\tiny, anchor=east] at (0,-1)    {1};
  \node[font=\tiny, anchor=east] at (0,-3.8)  {2};
\end{tikzpicture}
\caption{Diagramme de séquence --- Authentification}
\label{fig:auth_seq}
\end{figure}

\subsection{Assistant de Configuration de VM}

L'assistant est implémenté comme un formulaire React multi-étapes en sept étapes :
sélection du type d'instance, sélection du système d'exploitation, configuration du
stockage, configuration réseau, sélection de la région, révision et paiement. L'état
est géré globalement dans un contexte React \texttt{BuildVMContext}
(\texttt{frontend/BuildVMContext.tsx}), garantissant qu'aucune donnée n'est perdue
lors de la navigation entre les étapes.

\screenshotplaceholder{Assistant de configuration de VM en plusieurs étapes --- Étape 1 : Sélection du type d'instance}

Les images de VM sont répertoriées avec leurs UUID HCS directement dans le frontend.
Le catalogue distingue les images disponibles de celles marquées \texttt{needs\_id}
(UUID non encore attribué dans l'environnement HCS de production).

\subsection{Pipeline de Provisionnement de VM}

La figure~\ref{fig:vm_prov_seq} illustre la séquence de provisionnement de VM,
du paiement confirmé à la VM en cours d'exécution.

\begin{figure}[H]
\centering
\begin{tikzpicture}[
  box/.style={rectangle, draw, minimum width=2cm, minimum height=0.7cm,
              align=center, font=\scriptsize},
  arr/.style={-{Stealth}, thick},
]
  \node[box] (fe)  at (0,  0) {Frontend};
  \node[box] (be)  at (4,  0) {FastAPI};
  \node[box] (str) at (7.5,0) {Stripe};
  \node[box] (tf)  at (11, 0) {Terraform};
  \node[box] (hcs) at (14, 0) {HCS};

  \draw[dashed] (0,0) -- (0,-12);
  \draw[dashed] (4,0) -- (4,-12);
  \draw[dashed] (7.5,0) -- (7.5,-12);
  \draw[dashed] (11,0) -- (11,-12);
  \draw[dashed] (14,0) -- (14,-12);

  \draw[arr] (0,-1)   -- node[above,font=\tiny]{POST /api/payment/create-payment-intent} (4,-1);
  \draw[arr] (4,-1.8) -- node[above,font=\tiny]{confirm PaymentIntent} (7.5,-1.8);
  \draw[arr] (7.5,-2.6)-- node[above,font=\tiny]{succeeded} (4,-2.6);
  \draw[arr] (4,-3.4) -- node[above,font=\tiny]{\{session\_id\}} (0,-3.4);
  \node[font=\tiny, align=left, anchor=west] at (4.1,-4.0) {thread d'arrière-plan};
  \draw[arr] (4,-4.6) -- node[above,font=\tiny]{terraform init + apply} (11,-4.6);
  \draw[arr] (11,-5.4)-- node[above,font=\tiny]{HCS ECS créer + EIP} (14,-5.4);
  \draw[arr] (14,-6.2)-- node[above,font=\tiny]{VM en cours d'exécution} (11,-6.2);
  \draw[arr] (11,-7.0)-- node[above,font=\tiny]{ip\_publique, id\_vm} (4,-7.0);
  \draw[arr] (0,-7.8) -- node[above,font=\tiny]{GET /api/payment/progress/:session} (4,-7.8);
  \draw[arr] (4,-8.6) -- node[above,font=\tiny]{flux SSE (étapes)} (0,-8.6);
  \node[font=\tiny, align=left, anchor=west] at (4.1,-9.2) {arrière-plan : SSH + Netdata};
  \draw[arr] (4,-9.8) -- node[above,font=\tiny]{installer Netdata via SSH} (11,-9.8);
  \draw[arr] (4,-10.6)-- node[above,font=\tiny]{stocker VM + envoyer e-mail} (0,-10.6);
  \draw[arr] (4,-11.4)-- node[above,font=\tiny]{SSE : done=true} (0,-11.4);
\end{tikzpicture}
\caption{Diagramme de séquence --- Provisionnement de VM}
\label{fig:vm_prov_seq}
\end{figure}

\paragraph{Modèle Terraform.}
Chaque VM possède son propre répertoire d'état Terraform à
\texttt{terraform\_states/<id\_vm\_cloud>/}. Le modèle (\texttt{terraform\_test/main.tf})
déclare une \texttt{hcs\_ecs\_compute\_instance} avec un disque système, un EIP
(\texttt{hcs\_vpc\_eip}) et une association EIP. Le flavor est résolu par nom ou par
ID via une requête de source de données sur tous les flavors HCS disponibles.

\paragraph{Authentification SSH root via cloud-init.}
Les images Ubuntu 22.04 sur HCS imposent \texttt{PasswordAuthentication no} par
défaut. Le champ \texttt{user\_data} passe un script bash à \texttt{cloud-init} au
premier démarrage pour réactiver le SSH root par mot de passe, permettant au backend
de se connecter pour l'installation de Netdata.

\paragraph{Suivi de progression.}
La progression du provisionnement est suivie dans un dictionnaire Python en mémoire
\texttt{progress\_store} dans \texttt{backend/app/api/payment.py}. Chaque étape
(paiement confirmé, infrastructure provisionnée, supervision installée) est mise à
jour par le thread d'arrière-plan. Le frontend suit le flux SSE
\texttt{GET /api/payment/progress/:session\_id} et affiche un indicateur d'étapes.

\screenshotplaceholder{Indicateur de progression SSE en temps réel lors du provisionnement de VM}

\subsection{Tableau de Bord}

Le tableau de bord (\texttt{frontend/app/dashboard/}) liste toutes les VM appartenant
à l'utilisateur authentifié. Chaque carte VM affiche le statut, l'IP publique, le
flavor, la date de création et des boutons d'action pour le SSH, la supervision et
la suppression.

\screenshotplaceholder{Tableau de bord utilisateur --- liste des VM avec boutons d'action}

\subsection{Difficultés Rencontrées dans le Sprint 1}

\begin{itemize}
  \item \textbf{Nom vs.\ ID de flavor :} L'API HCS exige des ID de flavor (UUID), mais
        les utilisateurs et les modèles de l'assistant utilisent des noms de flavor
        (ex. : \texttt{c7.large.4}). Résolu en interrogeant
        \texttt{data.hcs\_ecs\_compute\_flavors.all} et en faisant correspondre par
        nom ou ID au moment du plan Terraform.
  \item \textbf{Type d'EIP incorrect :} Le type d'IP publique \texttt{5\_bgp} utilisé
        dans certains exemples Terraform n'est pas valide sur le déploiement HCS MESRS ;
        le type correct est \texttt{External-01}. Résolu en lisant les valeurs de
        l'interface portail HCS.
\end{itemize}

% =============================================================================
\section{Sprint 2 --- Recommandation IA, Terminal SSH et Supervision}
% =============================================================================

\subsection{Objectif du Sprint}

Ajouter le moteur de recommandation IA, un terminal SSH dans le navigateur, la
supervision automatisée Netdata et une marketplace de modèles de VM préconfigurés.

\subsection{Backlog du Sprint}

\begin{table}[H]
\centering
\caption{Backlog Sprint 2}
\label{tab:sprint2_backlog}
\begin{tabularx}{\linewidth}{llX}
\toprule
\textbf{ID} & \textbf{Statut} & \textbf{Description} \\
\midrule
US-08 & Terminé & Décrire sa charge de travail, recevoir une recommandation IA de VM \\
US-09 & Terminé & Parcourir les modèles de VM préconfigurés \\
US-10 & Terminé & Terminal SSH dans le navigateur \\
US-11 & Terminé & Tableau de bord de supervision Netdata par VM \\
\bottomrule
\end{tabularx}
\end{table}

\subsection{Moteur de Recommandation IA}

Le moteur IA s'exécute en tant que processus FastAPI distinct sur le port 8001 afin
d'isoler la latence d'inférence LLM de l'API principale. Il est composé de trois
modules :

\begin{enumerate}
  \item \textbf{\texttt{ollama\_client.py}} --- envoie la description textuelle de la
        charge de travail de l'utilisateur à l'API HTTP Ollama locale
        (\texttt{http://localhost:11434}) avec un prompt structuré demandant à
        Llama~3.1 de retourner un objet JSON.
  \item \textbf{\texttt{parser.py}} --- extrait et valide les champs JSON de la
        réponse du LLM : \texttt{workload} (intensité CPU/mémoire/GPU),
        \texttt{users} (nombre d'utilisateurs concurrents), \texttt{app\_type},
        \texttt{budget} et \texttt{confidence} (flottant entre 0 et 1).
  \item \textbf{\texttt{mapper.py}} --- mappe les champs analysés vers un flavor HCS
        et une spécification de configuration. L'intensité de la charge de travail
        (\texttt{low}/\texttt{medium}/\texttt{high}) sélectionne un niveau de flavor
        (ex. : \texttt{c7.large.4} pour les charges de calcul moyennes).
\end{enumerate}

La figure~\ref{fig:ai_flow} présente le flux de données de la recommandation IA.

\begin{figure}[H]
\centering
\begin{tikzpicture}[
  box/.style={rectangle, draw, rounded corners=3pt, minimum width=3cm,
              minimum height=0.9cm, align=center, font=\small},
  arr/.style={-{Stealth}, thick},
]
  \node[box, fill=green!10]  (user) at (0,0)  {Utilisateur\\(langage naturel)};
  \node[box, fill=orange!10] (be)   at (4,0)  {Backend\\FastAPI};
  \node[box, fill=purple!10] (ai)   at (8,0)  {Moteur IA};
  \node[box, fill=gray!10]   (oll)  at (8,-2) {Ollama\\llama3.1};
  \node[box, fill=blue!10]   (map)  at (12,0) {mapper.py\\flavor + specs};

  \draw[arr] (user) -- node[above,font=\scriptsize]{POST /api/ai/recommend} (be);
  \draw[arr] (be)   -- node[above,font=\scriptsize]{proxy} (ai);
  \draw[arr] (ai)   -- node[right,font=\scriptsize]{prompt} (oll);
  \draw[arr] (oll)  -- node[right,font=\scriptsize]{JSON} (ai);
  \draw[arr] (ai)   -- node[above,font=\scriptsize]{parse} (map);
  \draw[arr] (map)  to[bend right=30] node[below,font=\scriptsize]{recommandation} (be);
  \draw[arr] (be)   to[bend right=30] node[below,font=\scriptsize]{résultat} (user);
\end{tikzpicture}
\caption{Flux de données de la recommandation IA}
\label{fig:ai_flow}
\end{figure}

\screenshotplaceholder{Page de recommandation IA --- l'utilisateur saisit sa description de charge de travail, le système retourne la spécification de VM recommandée avec le score de confiance}

\subsection{Terminal SSH dans le Navigateur}

La fonctionnalité de terminal SSH offre un accès shell root à toute VM provisionnée
directement dans le navigateur, sans nécessiter l'installation d'un client SSH ni
la gestion de clés privées.

L'implémentation repose sur deux composants :

\begin{itemize}
  \item \textbf{Frontend :} La bibliothèque \texttt{xterm.js} rend un émulateur de
        terminal dans le navigateur et se connecte au backend via une connexion
        WebSocket.
  \item \textbf{Backend (\texttt{api/ssh.py}) :} Un endpoint WebSocket accepte la
        connexion, ouvre une session SSH Paramiko vers la VM cible (en utilisant l'IP
        publique de la VM et les identifiants root stockés), démarre un canal shell
        interactif et transfère les octets de manière bidirectionnelle entre le
        WebSocket et le canal SSH.
\end{itemize}

\screenshotplaceholder{Terminal SSH dans le navigateur affichant un shell root sur une VM HCS provisionnée}

\subsection{Supervision Netdata}

Après le provisionnement de la VM, un second thread d'arrière-plan se connecte à la
VM via SSH et installe l'agent Netdata :

\begin{lstlisting}[language=bash, caption={Commande d'installation de Netdata (exécutée via SSH)},
                   label={lst:netdata_install}]
bash <(curl -Ss https://my-netdata.io/kickstart.sh) --non-interactive
\end{lstlisting}

Le tableau de bord Netdata (port 19999) est ensuite accessible à
\texttt{http://<ip\_publique>:19999}. L'URL est stockée dans la colonne
\texttt{netdata\_url} de la table \texttt{virtual\_machines} et affichée sous forme
d'iframe sandboxé dans la page de supervision de la plateforme.

\screenshotplaceholder{Tableau de bord de supervision Netdata intégré dans la plateforme pour une VM en cours d'exécution}

\subsection{Marketplace de VM}

La page marketplace (\texttt{frontend/app/marketplace/}) affiche un catalogue de
modèles de VM préconfigurés (ex. : « Station Data Science », « Serveur Web »,
« Nœud ML Haute Mémoire ») avec les flavors recommandés, le système d'exploitation,
la taille du disque et le coût estimé. La sélection d'un modèle pré-remplit l'assistant
de création de VM.

% =============================================================================
\section{Sprint 3 --- Provisionnement de Clusters Kubernetes}
% =============================================================================

\subsection{Objectif du Sprint}

Concevoir et implémenter un pipeline complet de provisionnement de clusters Kubernetes
qui initialise des clusters multi-nœuds sur HCS avec kubeadm, dans les contraintes de
quota EIP de l'environnement cloud MESRS partagé.

\subsection{Backlog du Sprint}

\begin{table}[H]
\centering
\caption{Backlog Sprint 3}
\label{tab:sprint3_backlog}
\begin{tabularx}{\linewidth}{llX}
\toprule
\textbf{ID} & \textbf{Statut} & \textbf{Description} \\
\midrule
US-12 & Terminé & Configurer et payer un cluster Kubernetes \\
US-13 & Terminé & Provisionner le cluster K8s sur HCS avec kubeadm \\
US-14 & Terminé & Diffuser la progression du provisionnement K8s via SSE \\
US-15 & Terminé & Télécharger le kubeconfig depuis le tableau de bord \\
US-16 & Terminé & Supprimer un cluster K8s depuis le tableau de bord \\
\bottomrule
\end{tabularx}
\end{table}

\subsection{Architecture du Cluster}

Un cluster K8s sur HCS est composé d'un nœud maître (avec une EIP) et de $N$ nœuds
workers (IP privée uniquement). Cette topologie asymétrique a été imposée par le quota
EIP du projet HCS partagé MESRS : une seule EIP par cluster pouvait être allouée.

Le nœud maître remplit simultanément trois rôles :
\begin{enumerate}
  \item \textbf{Serveur API Kubernetes} (port 6443) --- accessible depuis l'extérieur
        via l'EIP du maître.
  \item \textbf{Hôte de rebond (bastion)} --- le backend se connecte aux nœuds workers
        via la session SSH du maître en utilisant l'API Paramiko
        \texttt{transport.open\_channel("direct-tcpip",...)}.
  \item \textbf{Passerelle NAT} --- le maître transfère les paquets à destination
        d'Internet provenant des nœuds workers, grâce à des règles
        \texttt{iptables} MASQUERADE.
\end{enumerate}

La figure~\ref{fig:k8s_arch} illustre l'architecture réseau du cluster.

\begin{figure}[H]
\centering
\begin{tikzpicture}[
  noeud/.style={rectangle, draw, rounded corners=3pt, minimum width=3cm,
               minimum height=1.2cm, align=center, font=\small},
  ext/.style={ellipse, draw, fill=blue!10, minimum width=2.5cm, align=center, font=\small},
  arr/.style={-{Stealth}, thick},
  darr/.style={-{Stealth}, thick, dashed},
]
  \node[ext]  (internet) at (0, 4)    {Internet};
  \node[noeud, fill=orange!15] (maitre)  at (0, 0)
    {Nœud Maître\\EIP : 193.95.x.x\\IP privée : 192.168.x.x};
  \node[noeud, fill=green!10]  (worker1) at (-3.5,-3) {Worker 1\\IP privée uniquement};
  \node[noeud, fill=green!10]  (workerN) at (3.5,-3)  {Worker N\\IP privée uniquement};
  \node[noeud, fill=gray!10]   (backend) at (-6, 0)   {Backend\\FastAPI};

  % Internet <-> Maître (EIP)
  \draw[arr] (maitre) -- node[right,font=\scriptsize]{EIP (ports 6443, 22)} (internet);

  % Maître <-> Workers (subnet VPC)
  \draw[arr] (maitre) -- node[left,font=\scriptsize]{Subnet VPC} (worker1);
  \draw[arr] (maitre) -- node[right,font=\scriptsize]{Subnet VPC} (workerN);

  % NAT
  \draw[darr] (worker1) to[bend left=20]
    node[left,font=\tiny]{passerelle par défaut = IP privée maître} (maitre);
  \draw[darr] (maitre) to[bend left=10]
    node[right,font=\tiny]{MASQUERADE} (internet);

  % Backend -> Maître (SSH direct)
  \draw[arr] (backend) -- node[above,font=\scriptsize]{SSH (clé RSA)} (maitre);

  % Backend -> Worker (jump)
  \draw[darr] (backend) to[bend right=20]
    node[below,font=\tiny]{SSH via hôte de rebond} (worker1);

  % Labels K8s
  \node[font=\tiny, anchor=south, align=center] at (0,-0.6) {kubeadm init\\Flannel CNI};
  \node[font=\tiny, anchor=north] at (-3.5,-4.2) {kubeadm join};
\end{tikzpicture}
\caption{Architecture réseau du cluster Kubernetes sur HCS}
\label{fig:k8s_arch}
\end{figure}

\subsection{Modèle Terraform du Cluster}

Le modèle Terraform du cluster (\texttt{terraform\_cluster/main.tf}) provisionne :
\begin{itemize}
  \item Une instance maître \texttt{hcs\_ecs\_compute\_instance} avec
        \texttt{source\_dest\_check = false} (requis pour le transfert de paquets NAT).
  \item Un EIP maître et une association EIP.
  \item $N$ instances workers \texttt{hcs\_ecs\_compute\_instance} (sans EIP).
  \item Tous les nœuds reçoivent une clé publique RSA générée dynamiquement via
        \texttt{cloud-config} dans \texttt{user\_data}, utilisée pour le SSH sans
        mot de passe.
\end{itemize}

\paragraph{Génération dynamique de paire de clés SSH.}
Une paire de clés RSA 2048 bits fraîche est générée par session de provisionnement
via Paramiko :

\begin{lstlisting}[language=python, caption={Génération de paire de clés RSA par cluster},
                   label={lst:keygen}]
def _generate_ssh_keypair():
    key = paramiko.RSAKey.generate(2048)
    public_key_line = f"ssh-rsa {key.get_base64()} k8s-cluster"
    return key, public_key_line
\end{lstlisting}

La clé publique est transmise à Terraform en tant que variable \texttt{ssh\_public\_key}
et injectée dans les nœuds maître et workers via le format \texttt{cloud-config} :

\begin{lstlisting}[language=bash, caption={user\_data cloud-init pour l'injection de clé SSH},
                   label={lst:cloudinit}]
#cloud-config
disable_root: false
ssh_authorized_keys:
  - ssh-rsa AAAA... k8s-cluster
\end{lstlisting}

Cela s'est avéré nécessaire car les images Ubuntu 22.04 sur HCS imposent
\texttt{PasswordAuthentication no} au niveau système, rejetant toute connexion SSH
par mot de passe quelles que soient les modifications de \texttt{sshd\_config}
tentées à l'exécution.

\subsection{Séquence d'Initialisation Kubernetes}

La fonction \texttt{install\_kubernetes()} dans
\texttt{backend/app/services/cluster\_service.py} exécute les étapes suivantes via SSH :

\begin{enumerate}
  \item \textbf{Attente de 120 s} pour que les VM terminent cloud-init et deviennent
        joignables.
  \item \textbf{Connexion au maître} via SSH direct en utilisant la clé RSA provisionnée.
  \item \textbf{Installation des prérequis sur le maître :} désactivation du swap,
        chargement des modules noyau (\texttt{overlay}, \texttt{br\_netfilter}),
        configuration des paramètres sysctl, installation de \texttt{containerd} et
        de \texttt{kubeadm/kubelet/kubectl} (K8s 1.29).
  \item \textbf{Activation du NAT sur le maître :} détection de l'interface réseau
        par défaut, activation du \texttt{ip\_forward} et ajout d'une règle
        MASQUERADE \texttt{iptables}.
  \item \textbf{Connexion aux workers via hôte de rebond} (canal \texttt{direct-tcpip}
        Paramiko à travers le maître).
  \item \textbf{Configuration des workers :} remplacement de la passerelle par défaut
        par l'IP privée du maître, remplacement du DNS (\texttt{/etc/resolv.conf})
        par les DNS Google (8.8.8.8), installation des prérequis.
  \item \textbf{kubeadm init sur le maître} avec le CIDR pod \texttt{10.244.0.0/16},
        en annonçant l'IP privée et en ajoutant l'IP publique comme
        \texttt{--apiserver-cert-extra-sans} pour l'accès externe.
  \item \textbf{Déploiement du CNI Flannel} via \texttt{kubectl apply}.
  \item \textbf{Génération de la commande join} et rattachement de tous les workers
        au cluster.
  \item \textbf{Récupération du kubeconfig} et réécriture de l'URL du serveur API
        de l'IP privée vers l'EIP du maître.
\end{enumerate}

\subsection{Solution Réseau : Maître comme Passerelle NAT}

\paragraph{Problème.}
Les nœuds workers n'ont pas d'EIP et donc aucune route internet directe. Sans accès
à Internet, \texttt{apt-get} ne peut pas télécharger les paquets K8s, et
\texttt{containerd} ne peut pas récupérer les images de conteneurs.

\paragraph{Investigation.}
Les tentatives initiales d'activation du transfert NAT avec uniquement des règles
MASQUERADE \texttt{iptables} sur le maître se sont révélées insuffisantes. Les paquets
envoyés par les workers via le maître étaient silencieusement abandonnés malgré la
règle iptables active.

\paragraph{Cause racine.}
L'hyperviseur HCS applique un contrôle \texttt{source\_dest\_check} sur chaque carte
réseau de VM : les paquets arrivant sur la carte du maître dont l'IP de destination
n'est pas celle du maître lui-même sont rejetés au niveau de l'hyperviseur, avant
d'atteindre la pile réseau Linux. Il s'agit d'un mécanisme anti-usurpation qui ne
peut pas être désactivé depuis l'intérieur de la VM.

\paragraph{Solution.}
L'attribut \texttt{source\_dest\_check = false} a été ajouté au bloc \texttt{network}
du maître dans Terraform, ce qui désactive le contrôle au niveau de l'hyperviseur via
l'API HCS :

\begin{lstlisting}[language=bash, caption={Terraform : désactivation de source\_dest\_check sur la carte réseau du maître},
                   label={lst:sdc}]
resource "hcs_ecs_compute_instance" "master" {
  ...
  network {
    uuid              = var.subnet_id
    source_dest_check = false  # autorise le transfert NAT
  }
}
\end{lstlisting}

Avec cette modification, le trafic internet des workers est correctement acheminé
via la règle MASQUERADE du maître vers Internet.

\subsection{Assistant et Tableau de Bord du Cluster}

L'assistant de configuration du cluster (\texttt{frontend/app/build-cluster/}) collecte
le nom du cluster, les flavors du maître et des workers, le nombre de workers, l'image
OS et la taille du disque. Après le paiement via Stripe, la progression du
provisionnement est diffusée vers le frontend via un endpoint SSE dédié
(\texttt{GET /api/clusters/progress/:session\_id}) avec six étapes labelisées.

\screenshotplaceholder{Assistant de provisionnement du cluster Kubernetes}

\screenshotplaceholder{Flux de progression SSE du cluster K8s --- les 6 étapes sont complètes}

Une fois le cluster en statut \texttt{running}, le tableau de bord l'affiche aux côtés
des VM ordinaires avec un bouton « Télécharger le kubeconfig ».

\screenshotplaceholder{Tableau de bord affichant un cluster Kubernetes en cours d'exécution avec l'IP du maître et le nombre de workers}

% ─────────────────────────────────────────────────────────────────────────────
\section{Conclusion}
% ─────────────────────────────────────────────────────────────────────────────

Trois sprints itératifs ont livré une plateforme cloud en libre-service complète :
le Sprint 1 a établi le cœur du provisionnement de VM ; le Sprint 2 a ajouté la
recommandation IA, l'accès SSH et la supervision ; le Sprint 3 a résolu le défi
non trivial du provisionnement de clusters Kubernetes dans les contraintes de quota
EIP de l'environnement partagé MESRS. Le schéma maître-comme-passerelle-NAT,
nécessitant un remplacement du \texttt{source\_dest\_check} au niveau Terraform,
a été la solution architecturale clé. Le chapitre suivant valide la plateforme au
regard de ses objectifs initiaux.
